\subsubsection{Vzájemná poloha přímek ve směrnicovém tvaru}
\label{subsubsec:Relative_position_slope_intercept_lines}

Ve směrnicovém tvaru je směr přímky určen přímo její směrnicí.
Vzájemnou polohu dvou nesvislých přímek proto můžeme určit porovnáním
jejich směrnic a úseků na ose $y$.

Uvažujme dvě přímky
\begin{align*}
    p &: y=k_1x+q_1,\\
    q &: y=k_2x+q_2.
\end{align*}

Jejich směr určují směrnice $k_1$ a $k_2$.

Jestliže
\[
    k_1\neq k_2,
\]
mají přímky různé směry, a proto jsou různoběžné. Jejich průsečík
získáme řešením rovnice
\[
    k_1x+q_1=k_2x+q_2.
\]

Odtud
\begin{equation}\label{eq:Slope_intercept_intersection_x}
    x_P=\frac{q_2-q_1}{k_1-k_2}.
\end{equation}

Druhou souřadnici průsečíku získáme dosazením hodnoty $x_P$ do rovnice
jedné z přímek.

Jestliže
\[
    k_1=k_2,
\]
mají přímky stejný směr. Potom rozlišujeme dva případy:
\begin{itemize}
    \item je-li $q_1\neq q_2$, jsou přímky rovnoběžné;
    \item je-li $q_1=q_2$, jsou přímky totožné.
\end{itemize}

Celý postup můžeme shrnout následovně:
\[
\begin{array}{c|c}
    \text{Podmínka}
    &
    \text{Vzájemná poloha přímek}
    \\
    \hline
    k_1\neq k_2
    &
    \text{přímky jsou různoběžné}
    \\[1mm]
    k_1=k_2,\ q_1\neq q_2
    &
    \text{přímky jsou rovnoběžné}
    \\[1mm]
    k_1=k_2,\ q_1=q_2
    &
    \text{přímky jsou totožné}
\end{array}
\]

\begin{example}\label{exa:Relative_position_slope_intercept_lines}
    Určete vzájemnou polohu přímek $p:y=2x-3$ a $q:y=-x+3$.

    Jejich směrnice jsou
    \[
        k_p=2,
        \qquad
        k_q=-1.
    \]
    Protože
    \[
        k_p\neq k_q,
    \]
    jsou přímky různoběžné.

    Souřadnici $x$ průsečíku získáme z rovnosti
    \begin{align*}
        2x-3&=-x+3,\\
        3x&=6,\\
        x&=2.
    \end{align*}

    Dosadíme do první rovnice:
    \[
        y=2\cdot2-3=1.
    \]

    Přímky se protínají v bodě
    \[
        P[2;1],
    \]
    a proto
    \[
        p\cap q=\{P[2;1]\}.
    \]
\end{example}

\paragraph{Kolmé přímky ve směrnicovém tvaru}
\label{subsubsec:Perpendicular_slope_intercept_lines}

Uvažujme dvě nesvislé přímky se směrnicemi $k_1$ a $k_2$. Jako jejich
směrové vektory můžeme podle vztahu
\eqref{eq:Direction_vector_from_slope} zvolit
\[
    \vec{u}_1=(1;k_1),
    \qquad
    \vec{u}_2=(1;k_2).
\]

Přímky jsou kolmé právě tehdy, když jsou kolmé jejich směrové vektory.
Jejich skalární součin proto musí být roven nule:
\begin{align*}
    \vec{u}_1\cdot\vec{u}_2
    &=
    (1;k_1)\cdot(1;k_2),\\
    0
    &=
    1+k_1k_2.
\end{align*}

Odtud dostáváme podmínku
\begin{equation}\label{eq:Perpendicular_slopes}
    k_1k_2=-1.
\end{equation}

Je-li $k_1\neq0$, můžeme směrnici kolmé přímky vyjádřit jako
\begin{equation}\label{eq:Perpendicular_slope}
    k_2=-\frac{1}{k_1}.
\end{equation}

Směrnice navzájem kolmých nesvislých přímek jsou tedy čísla
s opačnými znaménky a převrácenými absolutními hodnotami.

\begin{remark}\label{rem:Vertical_horizontal_perpendicular_lines}
    Vztah \eqref{eq:Perpendicular_slopes} lze použít pouze tehdy,
    když mají obě přímky definovanou směrnici. Svislá přímka nemá
    směrnici a je kolmá ke každé vodorovné přímce tvaru
    \[
        y=q.
    \]
\end{remark}

\begin{example}\label{exa:Perpendicular_line_slope_intercept}
    Určete směrnicový tvar přímky $q$, která prochází bodem $A[2;1]$ a je kolmá k přímce $p:y=2x+1$.

    Směrnice přímky $p$ je
    \[
        k_p=2.
    \]

    Směrnice kolmé přímky je
    \[
        k_q=-\frac{1}{2}.
    \]

    Použijeme bodově-směrnicový tvar:
    \begin{align*}
        y-1
        &=
        -\frac{1}{2}(x-2),\\
        y-1
        &=
        -\frac{1}{2}x+1,\\
        y
        &=
        -\frac{1}{2}x+2.
    \end{align*}

    Hledaná přímka má rovnici
    \[
        q:y=-\frac{1}{2}x+2.
    \]

    Kolmost můžeme ověřit součinem směrnic:
    \[
        2\cdot\left(-\frac{1}{2}\right)=-1.
    \]
\end{example}